\section{Proofs for the Solver, Certificates, and Recovery Laws}
\label{app:solver-recovery-proofs}

This appendix proves the algorithmic and recovery results in
\cref{sec:certified-solver,sec:multisecant-recovery}.
All gradients, exponential maps, and norms below
use the affine-invariant metric.

\subsection{Curvature and squared-distance comparison}

\begin{lemma}[AIRM curvature bound]
\label{lem:airm-curvature-bound}
Every sectional curvature of \(\SPD^d\), its determinant-one slice
\(\Pdet\), and every action fiber \(\Pfiber(B)\) lies in
\begin{equation}
  -\frac12\le \operatorname{sec}\le0.
  \label{eq:airm-curvature-bound}
\end{equation}
\end{lemma}

\begin{proof}
Congruence invariance reduces the calculation to the identity, where tangent
vectors are symmetric matrices and the metric is Frobenius.  The AIRM
curvature tensor is
\[
  R(X,Y)Z=-\frac14[[X,Y],Z],
\]
and hence
\[
  \operatorname{sec}(X,Y)
  =-\frac14\frac{\fro{[X,Y]}^2}
  {\fro X^2\fro Y^2-\ip X Y_F^2}\le0.
\]
To obtain the lower bound, replace \(Y\) by its component
\(\widetilde Y\) Frobenius-orthogonal to \(X\); this leaves the commutator
unchanged and turns the denominator into
\(\fro X^2\fro{\widetilde Y}^2\).  Orthogonally diagonalize
\(X=\operatorname{diag}(\lambda_i)\).  For symmetric
\(\widetilde Y=(y_{ij})\),
\[
\begin{aligned}
  \fro{[X,\widetilde Y]}^2
  &=2\sum_{i<j}(\lambda_i-\lambda_j)^2y_{ij}^2\\
  &\le4\fro X^2\sum_{i<j}y_{ij}^2
  \le2\fro X^2\fro{\widetilde Y}^2.
\end{aligned}
\]
Substitution gives \(\operatorname{sec}\ge-1/2\).  The determinant-one
slice and the action fibers are totally geodesic, so their intrinsic
curvatures are restrictions of the ambient curvature tensor.
\end{proof}

\begin{lemma}[Squared-distance Hessian comparison]
\label{lem:squared-distance-comparison}
Let \((\mathcal M,g)\) be a finite-dimensional Hadamard manifold with
\(-\kappa\le\operatorname{sec}\le0\).  For fixed \(z\), let
\(h(x)=\tfrac12d(z,x)^2\).  Then
\begin{equation}
  g_x\preceq\operatorname{Hess}h(x)
  \preceq\psi(\sqrt\kappa R)g_x
  \quad\text{whenever }d(z,x)\le R,
  \label{eq:squared-distance-comparison}
\end{equation}
where \(\psi(s)=s\coth s\) and \(\psi(0)=1\).  In particular, \(h\) is
globally \(1\)-strongly geodesically convex.
\end{lemma}

\begin{proof}
For \(x\ne z\), put \(r=d(z,x)\), let \(T\) be the terminal unit tangent
of the geodesic from \(z\) to \(x\), and decompose
\(v=aT+v_\perp\).  The radial Hessian formula gives
\[
  \operatorname{Hess}h(v,v)
  =a^2+r\operatorname{Hess}r(v_\perp,v_\perp).
\]
Let \(J\) be the Jacobi field with \(J(0)=0\) and
\(J(r)=v_\perp\).  Nonpositive curvature and Cauchy--Schwarz applied after
parallel transport yield
\[
  \operatorname{Hess}r(v_\perp,v_\perp)=I(J,J)
  \ge\int_0^r\norm{D_tJ}^2\,\dd t
  \ge r^{-1}\norm{v_\perp}^2.
\]
For the reverse bound, compare \(J\) with
\[
  V(t)=\frac{\sinh(\sqrt\kappa t)}
  {\sinh(\sqrt\kappa r)}P_tv_\perp,
\]
using \(V(t)=(t/r)P_tv_\perp\) when \(\kappa=0\).  The index lemma and the
curvature lower bound give
\[
  I(J,J)\le I(V,V)
  \le\sqrt\kappa\coth(\sqrt\kappa r)\norm{v_\perp}^2.
\]
Therefore
\[
  \norm v^2\le\operatorname{Hess}h(v,v)
  \le\psi(\sqrt\kappa r)\norm v^2.
\]
The formula extends continuously to \(x=z\), where
\(\operatorname{Hess}h=g_z\).  Monotonicity of \(\psi\) proves the stated
radius bound, and integration of the lower Hessian bound along geodesics
proves strong convexity.
\end{proof}

\begin{proof}[Proof of \cref{thm:canonical-fiber-engine}]
Coercivity and closedness give a minimizer, and strong convexity makes it
unique.  Write \(G_k=\operatorname{grad}F(x_k)\),
\(g_k=\norm{G_k}\), and
\[
  \gamma_k(t)=\operatorname{Exp}_{x_k}(-tG_k),
  \qquad 0\le t\le L_0^{-1}.
\]
If \(g_k=0\), strong convexity gives \(x_k=x_\star\).  Otherwise
\((F\circ\gamma_k)'(0)=-g_k^2<0\).  If this update segment left the current
sublevel, continuity would give a first return
\(t_\star\in(0,L_0^{-1}]\).  Before that return the Hessian is bounded by
\(L_0g\), so the integral Taylor formula would give
\[
\begin{aligned}
  F(\gamma_k(t_\star))
  &\le F(x_k)-t_\star g_k^2
       +\frac{L_0t_\star^2}{2}g_k^2\\
  &\le F(x_k)-\frac{t_\star}{2}g_k^2<F(x_k),
\end{aligned}
\]
contradicting the return.  Hence the complete update stays in the current
sublevel, and the same inequality at \(t=L_0^{-1}\) gives
\[
  F(x_{k+1})\le F(x_k)-\frac{g_k^2}{2L_0}.
\]

Let \(\Delta(x)=F(x)-F(x_\star)\),
\(d=d(x,x_\star)\), and
\(g=\norm{\operatorname{grad}F(x)}\).  Strong convexity at \(x\) and at
\(x_\star\) gives
\[
  \frac{\mu}{2}d^2
  \le\Delta(x)
  \le gd-\frac{\mu}{2}d^2
  \le\frac{g^2}{2\mu}.
\]
The first two expressions also imply \(d\le g/\mu\).  Thus
\(g^2\ge2\mu\Delta\), and the descent inequality yields
\[
  \Delta_{k+1}
  \le\left(1-\frac{\mu}{L_0}\right)\Delta_k.
\]
Iteration proves the value rate; the two preceding inequalities prove the
distance and posterior bounds.  Since
\(g_k^2/(2L_0)\le\Delta_k\), they also give the gradient rate.  Exact
inversion of \(q_{\mu,L}^k\) gives the stated complexity.  If
\(L_0=\mu\), the recurrence gives \(\Delta_1=0\).

For \(F(x)=\mu\|x\|^2/2\) on a Euclidean line, both posterior inequalities
and the \(L_0=\mu\) one-step statement hold with equality.
\end{proof}

For the fiber objective \(F_B(P)=\tfrac12\dai(I,P)^2\),
\cref{lem:airm-curvature-bound,lem:squared-distance-comparison} imply
\begin{equation}
  \operatorname{Hess}F_B\succeq g,
  \qquad
  \operatorname{Hess}F_B(P)\preceq
  \psi(D/\sqrt2)g_P
  \quad\text{if }\dai(I,P)\le D.
  \label{eq:fiber-hessian-bounds}
\end{equation}
The following action instance proves that the constant cannot be decreased
uniformly over dimensions and admissible actions.

\begin{proposition}[Action-family attainment of the radius majorant]
\label{prop:sharp-action-majorant}
Let \(d=3\), \(m=1\), and \(A=B=e_1\).  The action fiber is
\[
  \{\operatorname{diag}(1,\Sigma):\Sigma\in\mathscr P_2\},
\]
its canonical controller is \(I_3\), and for every \(D>0\) there are a point
at AIRM radius \(D\) and a unit fiber tangent \(v\) such that
\begin{equation}
  \operatorname{Hess}F_B(v,v)=\psi(D/\sqrt2).
  \label{eq:sharp-action-majorant-attainment}
\end{equation}
Hence \(\psi(D/\sqrt2)\) is the sharp dimension-uniform radius-only Hessian
majorant within the action family.  Along a radial geodesic in the same fiber,
both posterior inequalities in
\cref{eq:solver-distance-certificate,eq:solver-gap-certificate} hold with
equality.
\end{proposition}

\begin{proof}
Here \(M=1\), \(N=0\), and \(\rho=1\), so the displayed fiber and canonical
controller follow directly from \cref{eq:action-bundle-map}.  The
determinant-one surface \(\mathscr P_2\) is a two-dimensional space of
constant sectional curvature \(-1/2\), as equality holds in the commutator
bound of \cref{lem:airm-curvature-bound}.  On a constant-curvature
\(-\kappa\) plane, the transverse eigenvalue of the Hessian of
\(\tfrac12d(I,\cdot)^2\) at radius \(D\) is
\(D\sqrt\kappa\coth(D\sqrt\kappa)\).  Taking \(\kappa=1/2\) gives
\eqref{eq:sharp-action-majorant-attainment}.
\end{proof}

\subsection{Residual gradient and exact global convergence}

\begin{proof}[Proof of \cref{eq:residual-gradient-identity}]
On the ambient SPD manifold, for
\(\widetilde f(Q)=\tfrac12\dai(R_0,Q)^2\),
\[
  -\operatorname{grad}\widetilde f(Q)=\operatorname{Log}_Q(R_0).
\]
Whitening at \(Q\) turns this tangent into
\(\log(Q^{-1/2}R_0Q^{-1/2})\).  At
\(Q=\mathcal Q_\Sigma\) it is precisely \(\mathcal R(\Sigma)\).
The whitened tangent space of the hidden determinant-one block consists of
\[
  \begin{pmatrix}0&0\\0&Z\end{pmatrix},
  \qquad Z=Z^\top,\quad\tr Z=0.
\]
Since whitening converts the AIRM inner product to the Frobenius inner
product, orthogonal projection retains the \(22\) block and removes its
trace.  Thus the whitened reduced negative gradient is
\(W_B(\Sigma)=\mathcal R_{22}(\Sigma)^\circ\), which proves both identities
in \cref{eq:residual-gradient-identity}.
\end{proof}

\begin{proof}[Proof of \cref{thm:certified-action-solver}]
Set
\[
  f_k=f_B(\Sigma_k),\qquad
  g_k=\fro{W_B(\Sigma_k)},\qquad
  D_k=\dai(R_0,\mathcal Q_{\Sigma_k})=\sqrt{2f_k}.
\]
Trace freeness of
\(W_B(\Sigma_k)\) preserves determinant one, while the exponential preserves
positive definiteness.  The action chart consequently keeps every controller
on \(\Pfiber(B)\).

Let \(\gamma_k(t)\) be the negative-gradient geodesic, parameterized for
\(0\le t\le L_0^{-1}\), and put
\(\varphi_k=f_B\circ\gamma_k\).  If \(g_k=0\), the iterate is optimal.
Otherwise \(\varphi_k'(0)=-g_k^2<0\).  If the update segment left the
current objective sublevel, there would be a first return
\(t_\star\in(0,L_0^{-1}]\).  Before that return,
\(D(\gamma_k(t))\le D_k\le D_0\), so
\eqref{eq:fiber-hessian-bounds} bounds the Hessian by \(L_0g\).  The
integral Taylor formula would then imply
\[
  \varphi_k(t_\star)
  \le\varphi_k(0)-t_\star g_k^2
      +\frac{L_0t_\star^2}{2}g_k^2
  <\varphi_k(0),
\]
a contradiction.  Thus the complete update segment stays in the current
sublevel.  Applying the same Taylor bound at \(L_0^{-1}\) gives
\[
  f_{k+1}\le f_k-\frac1{L_0}g_k^2
  +\frac1{2L_0}g_k^2
  =f_k-\frac1{2L_0}g_k^2.
\]
Thus \(D_{k+1}\le D_k\), closing the induction and proving
\cref{eq:solver-descent} globally.

Let \(\Delta(\Sigma)=f_B(\Sigma)-f_B(\Sigma_\star)\).  Strong convexity
and Cauchy--Schwarz give, with
\(d=\dai(\Sigma,\Sigma_\star)\) and
\(g=\norm{\operatorname{grad}f_B(\Sigma)}\),
\begin{equation}
  \Delta(\Sigma)\le gd-\frac12d^2\le\frac12g^2.
  \label{eq:proof-pl}
\end{equation}
Hence \(g^2\ge2\Delta\), and the descent inequality yields
\[
  \Delta_{k+1}\le(1-L_0^{-1})\Delta_k=q_0\Delta_k.
\]
Since \(\Delta_0\le f_0=D_0^2/2\), this proves
\cref{eq:solver-value-rate}.  Strong convexity based at the minimizer gives
\(d^2\le2\Delta\), proving \cref{eq:solver-distance-rate}.  Combining
\(f_{k+1}\ge f_B(\Sigma_\star)\) with the descent inequality gives
\(g_k^2\le2L_0\Delta_k\), proving
\cref{eq:solver-residual-rate}.

The two strong-convexity inequalities also give
\(\Delta\ge d^2/2\).  Together with the first bound in
\eqref{eq:proof-pl}, this implies \(d\le g\), including the case \(d=0\),
while the second bound in \eqref{eq:proof-pl} gives
\(\Delta\le g^2/2\).  These are
\cref{eq:solver-distance-certificate,eq:solver-gap-certificate}.
Finally, for \(D_0>0\),
\(q_0^k=e^{-k\chi_0}\).  Solving the three rate bounds exactly for \(k\)
gives
\cref{eq:solver-value-complexity,eq:solver-distance-complexity,eq:solver-residual-complexity}.
The case \(D_0=0\) is immediate.
\end{proof}

\subsection{Active and inexact implementations}

\begin{proof}[Proof of \cref{cor:active-core-solver}]
Retain \(A=UC\) from \cref{sec:action-bundle}.  Choose \(V_a\) so that
\(E=[U,V_a]\) is an orthonormal basis of \(\mathcal W_A\).  Then
\[
  BC^{-1}=UM+V_aN_a,\qquad
  T_a=N_aM^{-1},\qquad
  L_a=\begin{pmatrix}I&0\\T_a&I\end{pmatrix},
  \qquad
  R_{0,a}=L_a^{-1}L_a^{-\top}.
\]
Complete \(E\) by \(Z\) when \(q>0\).  In the basis \([U,V_a,Z]\), the
full action coordinate satisfies
\[
  N=\begin{pmatrix}N_a\\0\end{pmatrix},
  \qquad
  L=L_a\oplus I_q,
  \qquad
  R_0=R_{0,a}\oplus I_q.
\]
An active-compatible initializer has the form
\[
  \Sigma=\Sigma_a\oplus\beta I_q,
  \qquad
  \det\Sigma_a\,\beta^q=1.
\]
The identity default corresponds to \(\Sigma_a=I_s\) and \(\beta=1\).
Then
\[
  \mathcal Q_a=\operatorname{diag}(M,\rho\Sigma_a),
  \qquad
  \mathcal Q_\Sigma
  =\mathcal Q_a\oplus(\rho\beta)I_q,
\]
and block functional calculus gives
\[
  \mathcal R_a
  =\log(\mathcal Q_a^{-1/2}R_{0,a}\mathcal Q_a^{-1/2}),
  \qquad
  \mathcal R(\Sigma)
  =\mathcal R_a\oplus r_cI_q,
  \qquad
  r_c=-\log(\rho\beta).
\]
The hidden trace-free projection subtracts the common mean
\[
  \tau=\frac{\tr\mathcal R_{22}^a+qr_c}{s+q}.
\]
Consequently,
\begin{equation}
  \begin{aligned}
    W_B(\Sigma)&=W_a\oplus w_cI_q,
    &W_a&=\mathcal R_{22}^a-\tau I_s,
    &w_c&=r_c-\tau,\\
    \fro{W_B(\Sigma)}^2&=\fro{W_a}^2+qw_c^2.
  \end{aligned}
\label{eq:active-residual-splitting}
\end{equation}
The full exponential update now splits exactly into
\begin{equation}
  \Sigma_a^+
  =\Sigma_a^{1/2}e^{W_a/L_0}\Sigma_a^{1/2},
  \qquad
  \beta^+=\beta e^{w_c/L_0}.
\label{eq:active-solver-update-proof}
\end{equation}
Since \(\tr W_a+qw_c=0\), this update preserves
\(\det\Sigma_a\,\beta^q=1\).  The action-bundle map then preserves
positive definiteness, determinant one, and \(PA=B\).  This proves
\eqref{eq:active-residual-splitting} and
\eqref{eq:active-solver-update-proof}; the
stopping certificate follows from
\cref{eq:solver-distance-certificate}.

A rank-revealing thin QR of the \(d\times2m\) input costs \(O(dm^2)\).
All remaining preprocessing acts on matrices of order at most \(r\), costing
\(O(r^3)\).  Each iteration uses one \(r\times r\) logarithm and inverse
square root and one \(s\times s\) exponential, hence \(O(r^3)\) arithmetic
and \(O(r^2)\) working memory.  Combining this cost with
\cref{eq:solver-residual-complexity} gives
\cref{eq:active-total-complexity}.  The representation
\[
  Px=cx+E(H-cI_r)E^\top x
\]
costs \(O(dr+r^2)\) per application.  If \(q=0\), every complement term is
absent and the same construction is the full solver.  If \(s=0\), the
determinant condition forces \(\beta=1\), leaving no variable.
\end{proof}

\begin{proof}[Proof of \cref{cor:inexact-action-solver}]
Put \(W_k=W_B(\Sigma_k)\), \(g_k=\fro{W_k}\), and
\(\widetilde W_k=W_k+E_k\).  First,
\cref{eq:computable-relative-error-test} and \(b_k\ge\fro{E_k}\) give
\[
  \fro{E_k}
  \le\frac{\delta}{1+\delta}
      \left(\fro{W_k}+\fro{E_k}\right),
\]
so rearrangement proves \cref{eq:relative-residual-error}.  Consequently,
\[
  (1-\delta)g_k\le\fro{\widetilde W_k}
  \le(1+\delta)g_k,
\qquad
  \ip{-W_k}{\widetilde W_k}_F
  \le-(1-\delta)g_k^2.
\]
If \(g_k=0\), the relative error forces a zero update.  Otherwise consider
the inexact update geodesic.  If it left the current objective sublevel,
there would be a first return \(t_\star\in(0,\eta]\).  Before that return,
the identity radius is at most \(D_k\le D_0\le\widehat D_0\), so the Hessian
is bounded by \(\widehat L_0g\).  The integral Taylor bound would give
\[
  f_B(\gamma_k(t_\star))
  \le f_k-t_\star(1-\delta)g_k^2
      +\frac{\widehat L_0t_\star^2}{2}(1+\delta)^2g_k^2
  <f_k,
\]
contradicting the return.  Hence the whole segment remains in the current
sublevel.  The descent lemma at \(t=\eta\) implies
\[
  f_{k+1}
  \le f_k-\eta(1-\delta)g_k^2
  +\frac{\widehat L_0\eta^2}{2}(1+\delta)^2g_k^2
  =f_k-\frac{c_\delta}{2\widehat L_0}g_k^2.
\]
Hence \(D_{k+1}\le D_k\).  Applying the global PL inequality from the exact
proof gives
\[
  \Delta_{k+1}\le\widehat q_\delta\Delta_k.
\]
Iteration and
\(\Delta_0\le D_0^2/2\le\widehat D_0^2/2\) prove
\cref{eq:inexact-value-rate}; strong convexity gives
\cref{eq:inexact-distance-rate}.  The descent bound and
\(f_{k+1}\ge f_B(\Sigma_\star)\) imply
\[
  g_k^2\le\frac{2\widehat L_0}{c_\delta}\Delta_k.
\]
Together with \(\fro{\widetilde W_k}\le(1+\delta)g_k\), the value rate, and
\(\sqrt{c_\delta}=(1-\delta)/(1+\delta)\), this proves
\cref{eq:inexact-residual-rate}.  Finally, the exact posterior certificates and
\(g_k\le\fro{\widetilde W_k}/(1-\delta)\) give
\cref{eq:inexact-certificates}.
\end{proof}

\subsection{Nested recovery}

\begin{proof}[Proof of \cref{thm:nested-multisecant}]
For every \(k\),
\[
  A_k^\top B_k
  =c_H(S^{(k)})^\top HS^{(k)}\succ0,
\]
so \(\mathscr C_k\) is nonempty, closed, and totally geodesic by
\cref{thm:action-bundle}.  Also
\(P_\star A_k=c_HH^{-1}HS^{(k)}=B_k\), so
\(P_\star\in\mathscr C_k\).  Nested column spaces imply
\(\mathscr C_{k+1}\subset\mathscr C_k\).  Metric projection onto a
nonempty closed geodesically convex subset of a Hadamard manifold exists
uniquely, proving that every update is well defined and retains every earlier
action.

For completeness, if \(y=\operatorname{Proj}_C(x)\) and \(z\in C\), first
variation along the geodesic from \(y\) to \(z\) gives
\(\ip{\operatorname{Log}_y(x)}{\operatorname{Log}_y(z)}\le0\).  Since the
exponential map is distance expanding,
\[
\begin{aligned}
  d(x,z)^2
  &\ge\norm{\operatorname{Log}_y(x)-\operatorname{Log}_y(z)}^2\\
  &\ge d(x,y)^2+d(y,z)^2.
\end{aligned}
\]
Apply this with \(x=P_{k-1}\), \(y=P_k\), \(z=P_\star\), and
\(C=\mathscr C_k\) to obtain
\cref{eq:multisecant-pythagorean}.  Summation telescopes to
\cref{eq:multisecant-energy}; Cauchy--Schwarz then gives
\cref{eq:multisecant-length}.

If \(r_K=d-1\), the bundle fiber is diffeomorphic to
\(\mathscr P_{d-r_K}=\mathscr P_1=\{1\}\), so
\(\mathscr C_K=\{P_\star\}\).  If \(r\le d-2\), then
\(n=d-r\ge2\) and the fiber is diffeomorphic to \(\mathscr P_n\), whose
dimension \(n(n+1)/2-1\) is positive; it therefore contains infinitely many
completions.  This proves the sharp threshold and the stated block count.

Finally, for the \(k\)-th projection set
\[
  \widetilde P=P_{k-1}^{-1/2}PP_{k-1}^{-1/2},\quad
  \widetilde A_k=P_{k-1}^{1/2}A_k,\quad
  \widetilde B_k=P_{k-1}^{-1/2}B_k.
\]
Then \(PA_k=B_k\) is equivalent to
\(\widetilde P\widetilde A_k=\widetilde B_k\), determinant one is
preserved, and congruence invariance turns the objective into
\(\dai(I,\widetilde P)\).  Thus every outer update is an identity-reference
action problem.  The exact projection is the limit covered by
\cref{thm:certified-action-solver}; \cref{cor:inexact-action-solver} certifies
finite inner approximations but is not used in the exact Pythagorean identities.
\end{proof}

\begin{proof}[Proof of \cref{thm:approximate-multisecant}]
Fix a round and abbreviate
\[
  x=\widetilde P_{k-1},\qquad y=Q_k,\qquad
  \widetilde y=\widetilde P_k,\qquad z=P_\star,
\]
and let \(a=\dai(x,y)\), \(b=\dai(y,z)\),
\(e=\dai(y,\widetilde y)\), and \(D=\dai(x,z)\).
Because \(y\) is the projection of \(x\) onto \(\mathscr C_k\) and
\(z\in\mathscr C_k\), the CAT(0) projection inequality gives
\[
  a^2+b^2\le D^2.
\]
The triangle inequality and \(e\le\varepsilon_k\) give
\[
\begin{aligned}
  \dai(x,\widetilde y)^2+\dai(\widetilde y,z)^2
  &\le(a+e)^2+(b+e)^2\\
  &\le D^2+2e(a+b)+2e^2\\
  &\le D^2+2\sqrt2\,eD+2e^2,
\end{aligned}
\]
where the last line uses
\(a+b\le\sqrt{2(a^2+b^2)}\le\sqrt2D\).  This proves
\cref{eq:approximate-multisecant-pythagorean}.  Also
\[
  D_k\le b+e\le D_{k-1}+\varepsilon_k,
\]
and induction proves \cref{eq:approximate-multisecant-distance}.  Rearranging
\cref{eq:approximate-multisecant-pythagorean} and summing proves
\cref{eq:approximate-multisecant-energy}.  Nested feasibility preserves the
earlier actions.  The stated observable bounds are the posterior certificates
after the congruence in the preceding proof.  Finally, \(r_K=d-1\) makes the
final fiber a singleton by \cref{thm:nested-multisecant}.
\end{proof}
